\chapter{Attention Mechanisms, from Dot Products to FlashAttention}\label{ch:attention}

Attention is the one operation in a transformer that mixes information \emph{across} positions.
Everything else---embeddings, layer norms, MLPs---acts on each token independently. It is also the
one operation whose cost grows quadratically with sequence length, and so the one that dominates
long-context training, shapes the KV cache that dominates serving memory, and motivated the most
influential systems paper of the transformer era, FlashAttention.

This chapter builds attention up in small steps: a weight-free version that exposes the core
``weighted average of the inputs'' idea, then trainable queries, keys and values, then causal
masking and dropout, then multi-head attention in two implementations. With the algorithm in
hand, we analyse its FLOPs and memory traffic, derive the online softmax, and implement an
IO-aware tiled attention whose forward and backward passes match the naive version to machine
precision without ever forming the $T\times T$ attention matrix.

\begin{figure}[t]
\centering
\begin{tikzpicture}[
  st/.style={draw=sysblue,fill=sysbluebg,rounded corners=2pt,minimum height=11mm,text width=22mm,align=center,font=\small},
  sy/.style={st,draw=napkin,fill=napkinbg},
  num/.style={circle,fill=sysblue,text=white,font=\scriptsize\bfseries,inner sep=1.5pt},
  arr/.style={-{Stealth[length=2.2mm]},thick,sysblue!70}]
\node[st] (a) {Simplified\\self-attention};
\node[st,right=5mm of a] (b) {Self-attention\\with $\mat{W}_q,\mat{W}_k,\mat{W}_v$};
\node[st,right=5mm of b] (c) {Causal mask\\+ dropout};
\node[st,right=5mm of c] (d) {Multi-head\\attention};
\node[sy,below=9mm of b] (e) {Cost model:\\FLOPs \& bytes};
\node[sy,right=5mm of e] (f) {Online softmax\\+ tiling};
\node[sy,right=5mm of f] (g) {Variants:\\GQA, windows, \dots};
\foreach \n/\l in {a/1,b/2,c/3,d/4} \node[num] at (\n.north west) {\l};
\foreach \n/\l in {e/5,f/6,g/7} \node[num,fill=napkin] at (\n.north west) {\l};
\draw[arr] (a)--(b); \draw[arr] (b)--(c); \draw[arr] (c)--(d);
\draw[arr,napkin] (d.south) |- ++(0,-4mm) -| (e.north);
\draw[arr,napkin] (e)--(f); \draw[arr,napkin] (f)--(g);
\node[font=\scriptsize,text=srcgray,anchor=east] at ([xshift=-3mm]e.west) {systems half $\rightarrow$};
\end{tikzpicture}
\caption{Roadmap: steps 1--4 build the attention used in GPT-2 (blue); steps 5--7 analyse and
optimise it (amber).}
\label{fig:attention:roadmap}
\end{figure}

% =====================================================================
\section{The problem with modeling long sequences}\label{sec:attention:long}

Consider translating a sentence. Word-by-word substitution fails immediately: languages reorder
words, split one word into several, and resolve meanings from context far away in the sentence.
A translation model must read the whole source, keep what matters, and consult the relevant parts
as it writes each output word.

Before transformers, the standard architecture was the \emph{encoder--decoder recurrent neural
network} (Sutskever et al., 2014; Cho et al., 2014). The encoder RNN reads the source tokens one at
a time, updating a hidden state $\vect{h}_t = f(\vect{h}_{t-1},\vect{x}_t)$. Its final hidden
state---a single fixed-size vector---is handed to the decoder, which generates the output one token
at a time from that vector and its own running state (\cref{fig:attention:rnn}, left).

The weakness is the bottleneck: the entire meaning of an arbitrarily long input must be squeezed
through one vector of a few hundred numbers, and the decoder has no way to look back at individual
input positions. Empirically, translation quality fell off sharply for long sentences. A second,
systems-flavoured weakness is that the recurrence is \emph{sequential}: $\vect{h}_t$ cannot be
computed before $\vect{h}_{t-1}$, so training cannot parallelise across the time dimension of a
sequence, leaving most of a GPU idle.

\begin{figure}[t]
\centering
\begin{tikzpicture}[
  cell/.style={draw=fcompute,fill=sysbluebg,minimum size=7mm,rounded corners=1pt,font=\scriptsize},
  dcell/.style={draw=fdata,fill=keygreenbg,minimum size=7mm,rounded corners=1pt,font=\scriptsize},
  inp/.style={font=\scriptsize\itshape},
  arr/.style={-{Stealth},thick}]
% left: RNN bottleneck
\foreach \i/\w in {1/the,2/cache,3/is,4/full} {
  \node[cell] (e\i) at ({\i*1.0},0) {$\vect{h}_{\i}$};
  \node[inp] at ({\i*1.0},-0.75) {\w};
  \draw[arr,srcgray] ({\i*1.0},-0.6) -- (e\i);
}
\foreach \i [evaluate=\i as \j using int(\i+1)] in {1,2,3} \draw[arr] (e\i) -- (e\j);
\node[draw=warnred,fill=warnredbg,circle,minimum size=7mm,font=\scriptsize] (z) at (5.4,0) {$\vect{c}$};
\draw[arr,warnred] (e4) -- (z);
\foreach \i/\w in {1/der,2/Cache,3/ist,4/voll} {
  \node[dcell] (d\i) at ({\i*1.0},1.9) {$\vect{s}_{\i}$};
  \node[inp] at ({\i*1.0},2.65) {\w};
  \draw[arr,srcgray] (d\i) -- ({\i*1.0},2.5);
}
\foreach \i [evaluate=\i as \j using int(\i+1)] in {1,2,3} \draw[arr] (d\i) -- (d\j);
\draw[arr,warnred] (z) to[out=90,in=0] ([yshift=0mm]d4.east);
\draw[arr,warnred,dashed] (z) to[out=120,in=-30] (d1.south east);
\node[font=\scriptsize,text=warnred,align=center] at (5.4,-0.75) {single\\context vector};
\node[font=\small\bfseries] at (2.6,3.3) {RNN encoder--decoder};
% right: attention
\begin{scope}[xshift=7.3cm]
\foreach \i/\w in {1/the,2/cache,3/is,4/full} {
  \node[cell] (f\i) at ({\i*1.0},0) {$\vect{h}_{\i}$};
  \node[inp] at ({\i*1.0},-0.75) {\w};
}
\node[dcell] (g) at (2.5,1.9) {$\vect{s}_{2}$};
\node[inp] at (2.5,2.65) {Cache};
\foreach \i/\o in {1/25,2/90,3/40,4/20} \draw[arr,fcomm,line width=0.2+\o*0.025 mm,opacity=0.85] (f\i) -- (g);
\node[font=\scriptsize,text=fcomm,align=left,anchor=west] at (4.2,1.0) {weights\\$\alpha_{2,j}$};
\node[font=\small\bfseries] at (2.5,3.3) {decoder with attention};
\end{scope}
\end{tikzpicture}
\caption{Left: the encoder--decoder RNN must compress the whole input into one vector
$\vect{c}$. Right: with attention, each decoder step forms its own weighted combination of
\emph{all} encoder states; line width shows how strongly each input is used to produce
``Cache''.}
\label{fig:attention:rnn}
\end{figure}

% =====================================================================
\section{Capturing data dependencies with attention}\label{sec:attention:deps}

Bahdanau et al.\ (2015) removed the bottleneck by letting the decoder, at every output step,
compute a fresh set of weights over \emph{all} encoder hidden states and use their weighted sum as
context (\cref{fig:attention:rnn}, right). The weights are produced by a small learned scoring
function and normalised with a softmax, so the decoder can ``attend'' to whichever input positions
are relevant to the word it is producing.

Three years later, Vaswani et al.\ (2017) asked what happens if one drops the recurrence entirely
and builds the whole network out of attention. The resulting \emph{transformer} uses
\emph{self-attention}: a sequence attends to \emph{itself}, so every position's new representation
is a learned, data-dependent mixture of every (allowed) position's old representation. Without
recurrence, all positions of a sequence are processed in parallel by large matrix multiplications,
which is precisely what GPUs are built for. That property, more than any modelling advantage, is
why transformers displaced RNNs.

\begin{keyidea}[Self-attention in one sentence]
For each position $i$, compute a probability distribution $\alpha_{i,\cdot}$ over positions,
based on how well $i$'s ``question'' matches each position's ``label'', and output the
$\alpha$-weighted average of the positions' ``contents''. Queries, keys and values are those
three roles.
\end{keyidea}

% =====================================================================
\section{Self-attention without trainable weights}\label{sec:attention:simple}

To isolate the mechanism, we first drop all learned parameters. Take a six-token input whose
3-dimensional embeddings we simply write down (real models use $d$ in the hundreds or thousands
and learn the embeddings, \cref{ch:text}):

\begin{lstlisting}[caption={The running example: six tokens with toy 3-d embeddings.},label={lst:attention:inputs}]
inputs = torch.tensor(
  [[0.40, 0.10, 0.80],   # fused     (x1)
   [0.60, 0.90, 0.20],   # kernels   (x2)
   [0.50, 0.80, 0.70],   # keep      (x3)
   [0.10, 0.30, 0.90],   # data      (x4)
   [0.90, 0.20, 0.10],   # on        (x5)
   [0.20, 0.70, 0.40]])  # chip      (x6)
\end{lstlisting}

\subsection{One query at a time}
Pick the third token, ``keep'', as the \emph{query}, $\vect{q}=\vect{x}^{(3)}$. Self-attention
computes its \emph{context vector} $\vect{z}^{(3)}$ in three steps.
\begin{enumerate}
  \item \textbf{Scores.} Measure the similarity of the query to every input with a dot product:
  $\omega_{3j} = \vect{x}^{(3)}\cdot\vect{x}^{(j)}$. A dot product is large when two vectors point the
  same way and have large norms, so it serves as an un-normalised relevance score. For our numbers:
  $\vect{\omega}_3 = (0.84,\ 1.16,\ 1.38,\ 0.92,\ 0.68,\ 0.94)$.
  \item \textbf{Weights.} Normalise the scores into a probability distribution with a softmax,
  $\alpha_{3j} = e^{\omega_{3j}}/\sum_k e^{\omega_{3k}}$, giving
  $\vect{\alpha}_3 = (0.140, 0.193, 0.241, 0.152, 0.120, 0.155)$. Softmax keeps every weight positive,
  makes them sum to one, and exaggerates differences (it is a smooth arg-max).
  \item \textbf{Context.} Average the inputs with those weights:
  $\vect{z}^{(3)} = \sum_j \alpha_{3j}\vect{x}^{(j)} = (0.446,\ 0.558,\ 0.530)$.
\end{enumerate}
\Cref{fig:attention:onequery} draws the computation. Note that the query attends most strongly to
itself (its dot product with itself is its squared norm), a pattern that trainable projections
will be free to change.

\begin{figure}[t]
\centering
\begin{tikzpicture}[
  vec/.style={draw=fcompute,fill=sysbluebg,minimum width=15mm,minimum height=6mm,font=\scriptsize},
  qv/.style={vec,draw=warnred,fill=warnredbg},
  w/.style={font=\scriptsize,text=fcomm},
  arr/.style={-{Stealth},thick}]
\foreach \w/\v [count=\i] in {fused/{.4,.1,.8},kernels/{.6,.9,.2},keep/{.5,.8,.7},data/{.1,.3,.9},on/{.9,.2,.1},chip/{.2,.7,.4}} {
  \node[vec] (x\i) at ({(\i-1)*2.15},0) {$(\v)$};
  \node[font=\scriptsize\itshape] at ({(\i-1)*2.15},0.6) {\w};
}
\node[qv] at (x3) {$(.5,.8,.7)$};
\foreach \s/\a [count=\i] in {0.84/0.140,1.16/0.193,1.38/0.241,0.92/0.152,0.68/0.120,0.94/0.155} {
  \node[font=\scriptsize] (s\i) at ({(\i-1)*2.15},-1.0) {$\omega_{3\i}{=}\s$};
  \node[w] (a\i) at ({(\i-1)*2.15},-1.6) {$\alpha_{3\i}{=}\a$};
  \draw[arr,srcgray] (x\i) -- (s\i);
}
\draw[decorate,decoration={brace,amplitude=4pt,mirror},thick] (-0.8,-1.95) -- (11.55,-1.95)
  node[midway,below=5pt,font=\scriptsize] {softmax over the six scores};
\node[draw=fdata,fill=keygreenbg,minimum width=30mm,minimum height=7mm,font=\small] (z) at (5.4,-3.3)
  {$\vect{z}^{(3)}=\sum_j\alpha_{3j}\vect{x}^{(j)}=(.446,.558,.530)$};
\foreach \i in {1,...,6} \draw[arr,fcomm!70] (a\i.south) -- (z.north);
\node[font=\scriptsize,text=warnred,anchor=west] at (11.6,0) {query $=\vect{x}^{(3)}$};
\end{tikzpicture}
\caption{Weight-free self-attention for the query ``keep'': dot-product scores against every
input, softmax to weights, weighted sum to the context vector.}
\label{fig:attention:onequery}
\end{figure}

\begin{pitfall}[Never write softmax as \texttt{exp(x)/exp(x).sum()}]
A score of 100 overflows \code{exp} in fp16 (max $\approx 65504$) and one of 1000 overflows fp32.
Since $\softmax(\vect{x}) = \softmax(\vect{x}-c)$ for any constant $c$, numerically stable
implementations subtract the row maximum first. \code{torch.softmax} does this; so must every
hand-written kernel. The running maximum is also the key to the online softmax of
\cref{sec:attention:flash}.
\end{pitfall}

\subsection{All queries at once}
Repeating the procedure for every query is just three matrix operations on the input matrix
$\mat{X}\in\R^{T\times d}$:
\begin{equation}
\mat{\Omega} = \mat{X}\mat{X}^\top \in\R^{T\times T},\qquad
\mat{A} = \softmax_{\text{row}}(\mat{\Omega}),\qquad
\mat{Z} = \mat{A}\mat{X}\in\R^{T\times d}.
\label{eq:attention:simple}
\end{equation}
Row $i$ of $\mat{A}$ is the distribution of query $i$ over all keys (\cref{fig:attention:heat}).
In PyTorch it is three lines, and replacing Python loops with these matrix operations is the
first and most important performance optimisation in this chapter: one $T\times d$ by $d\times T$
GEMM instead of $T^2$ tiny dot products.

\begin{lstlisting}[caption={Weight-free self-attention for all tokens (Lab 4, Part A).},label={lst:attention:simple}]
def simple_attention(x):              # x: (T, d)
    scores = x @ x.T                  # (T, T) all pairwise dot products
    weights = torch.softmax(scores, dim=-1)   # each row sums to 1
    return weights @ x, weights       # (T, d) context vectors
\end{lstlisting}

\newcommand{\attnheatcell}[4]{% x, y, value, scale
  \pgfmathsetmacro{\shade}{min(100,#3*#4)}
  \fill[fcompute!\shade] (#1,-#2) rectangle ++(1,-1);
  \node[font=\tiny] at (#1+0.5,-#2-0.5) {#3};}
\begin{figure}[t]
\centering
\begin{tikzpicture}[x=8mm,y=8mm]
\foreach \w [count=\i from 0] in {fused,kernels,keep,data,on,chip} {
  \node[font=\scriptsize,anchor=east] at (-0.1,-\i-0.5) {\w};
  \node[font=\scriptsize,anchor=west,rotate=60] at (\i+0.3,0.1) {\w};
}
\foreach \r/\vals in {0/{.19,.14,.20,.19,.14,.14},1/{.11,.24,.22,.12,.15,.16},2/{.14,.19,.24,.15,.12,.15},
                      3/{.18,.14,.21,.21,.11,.15},4/{.15,.20,.18,.12,.22,.13},5/{.14,.20,.22,.15,.12,.17}} {
  \foreach \v [count=\c from 0] in \vals { \attnheatcell{\c}{\r}{\v}{280} }
}
\draw[srcgray] (0,0) rectangle (6,-6);
\node[font=\small,anchor=north] at (3,-6.2) {(a) weight-free, all keys};
\begin{scope}[xshift=64mm]
\foreach \w [count=\i from 0] in {fused,kernels,keep,data,on,chip} {
  \node[font=\scriptsize,anchor=west,rotate=60] at (\i+0.3,0.1) {\w};
}
\foreach \r/\vals in {0/{1.0},1/{.33,.67},2/{.24,.34,.42},3/{.25,.19,.28,.28},4/{.17,.23,.21,.14,.25},5/{.14,.20,.22,.15,.12,.17}} {
  \foreach \v [count=\c from 0] in \vals { \attnheatcell{\c}{\r}{\v}{140} }
}
\foreach \r in {0,...,4} { \foreach \c in {\r,...,4} {
  \pgfmathtruncatemacro{\cc}{\c+1}
  \fill[pattern=north east lines,pattern color=warnred!50] (\cc,-\r) rectangle ++(1,-1); } }
\draw[srcgray] (0,0) rectangle (6,-6);
\node[font=\small,anchor=north] at (3,-6.2) {(b) causal (\cref{sec:attention:causal})};
\end{scope}
\end{tikzpicture}
\caption{Attention weight matrices $\mat{A}$ for the running example (rows: queries, columns:
keys; numbers computed with \cref{lst:attention:simple}). Each row sums to 1. In (b) the hatched
future positions are masked and each row is renormalised over the visible prefix.}
\label{fig:attention:heat}
\end{figure}

% =====================================================================
\section{Self-attention with trainable weights}\label{sec:attention:trainable}

The weight-free version has no parameters to learn: what a token attends to is fixed by raw
embedding similarity, and what it retrieves is the raw embedding itself. \emph{Scaled dot-product
attention} (Vaswani et al., 2017) introduces three learned projection matrices
$\mat{W}_q,\mat{W}_k\in\R^{d_{\text{in}}\times d_k}$ and $\mat{W}_v\in\R^{d_{\text{in}}\times d_v}$ that
give each token three separate roles:
\begin{itemize}
  \item \textbf{query} $\vect{q}_i=\vect{x}_i\mat{W}_q$ --- what token $i$ is looking for;
  \item \textbf{key} $\vect{k}_j=\vect{x}_j\mat{W}_k$ --- what token $j$ advertises about itself;
  \item \textbf{value} $\vect{v}_j=\vect{x}_j\mat{W}_v$ --- what token $j$ contributes if attended to.
\end{itemize}
The terminology comes from databases and key--value stores: a query is matched against keys, and
the matching entries' values are returned---except that here matching is soft and every value is
returned with some weight.

\subsection{Computing the attention weights step by step}
For the full input,
\begin{equation}
\mat{Q}=\mat{X}\mat{W}_q,\quad \mat{K}=\mat{X}\mat{W}_k,\quad \mat{V}=\mat{X}\mat{W}_v,\qquad
\operatorname{Attention}(\mat{Q},\mat{K},\mat{V}) = \softmax\!\Big(\frac{\mat{Q}\mat{K}^\top}{\sqrt{d_k}}\Big)\mat{V}.
\label{eq:attention:sdpa}
\end{equation}
\Cref{fig:attention:qkv} traces the shapes for our example with $d_{\text{in}}=3$ and
$d_k=d_v=2$ (deliberately different so the shapes cannot be confused; in GPT models
$d_{\text{in}}=d_{\text{out}}$).

\begin{figure}[t]
\centering
\begin{tikzpicture}[
  m/.style={draw,minimum height=#1,align=center,font=\scriptsize,inner sep=1pt},
  arr/.style={-{Stealth},thick}]
\node[m=16mm,minimum width=9mm,fill=white] (X) {$\mat{X}$\\$6{\times}3$};
\node[m=9mm,minimum width=6mm,fill=warnredbg,draw=warnred,right=12mm of X,yshift=13mm] (Wq) {$\mat{W}_q$\\$3{\times}2$};
\node[m=9mm,minimum width=6mm,fill=sysbluebg,draw=fcompute,right=12mm of X] (Wk) {$\mat{W}_k$\\$3{\times}2$};
\node[m=9mm,minimum width=6mm,fill=keygreenbg,draw=fdata,right=12mm of X,yshift=-13mm] (Wv) {$\mat{W}_v$\\$3{\times}2$};
\node[m=16mm,minimum width=6mm,fill=warnredbg,draw=warnred,right=10mm of Wq] (Q) {$\mat{Q}$\\$6{\times}2$};
\node[m=16mm,minimum width=6mm,fill=sysbluebg,draw=fcompute,right=10mm of Wk] (K) {$\mat{K}$\\$6{\times}2$};
\node[m=16mm,minimum width=6mm,fill=keygreenbg,draw=fdata,right=10mm of Wv] (V) {$\mat{V}$\\$6{\times}2$};
\foreach \w in {Wq,Wk,Wv} \draw[arr,srcgray] (X.east) -- (\w.west);
\draw[arr] (Wq)--(Q); \draw[arr] (Wk)--(K); \draw[arr] (Wv)--(V);
\node[m=16mm,minimum width=16mm,fill=labpurplebg,draw=fcomm,right=14mm of K,yshift=6mm] (S) {$\mat{Q}\mat{K}^\top/\sqrt{2}$\\$6{\times}6$};
\draw[arr] (Q.east) -- (S.west); \draw[arr] (K.east) -- (S.west);
\node[m=16mm,minimum width=16mm,fill=labpurplebg,draw=fcomm,right=8mm of S] (A) {$\mat{A}=\softmax(\cdot)$\\$6{\times}6$};
\draw[arr] (S)--(A);
\node[m=16mm,minimum width=8mm,fill=napkinbg,draw=fmemory,below=8mm of A] (Z) {$\mat{Z}=\mat{A}\mat{V}$\\$6{\times}2$};
\draw[arr] (A)--(Z); \draw[arr] (V.east) -- (Z.west);
\end{tikzpicture}
\caption{Shapes in single-head scaled dot-product attention for the running example
($T=6$, $d_{\text{in}}=3$, $d_k=d_v=2$).}
\label{fig:attention:qkv}
\end{figure}

\paragraph{Why divide by $\sqrt{d_k}$?} Suppose the entries of $\vect{q}$ and $\vect{k}$ are independent
with mean 0 and variance 1. Then $\vect{q}\cdot\vect{k}=\sum_{m=1}^{d_k}q_mk_m$ has mean 0 and variance
$d_k$. With $d_k=128$, raw scores have a standard deviation of about 11; a softmax over such scores
is nearly one-hot, and its gradient (which is proportional to $\alpha_i(1-\alpha_i)$ on the diagonal)
nearly vanishes. Dividing by $\sqrt{d_k}$ restores unit variance at initialisation and keeps the
softmax in a regime where gradients flow. (Some newer models instead normalise $\vect{q}$ and
$\vect{k}$ directly---the QK-norm of Qwen3, \cref{ch:modern}.)

\subsection{A compact self-attention class}
Wrapping the computation in an \code{nn.Module} gives us trainable parameters, device movement and
state-dict serialisation for free:

\begin{lstlisting}[caption={Single-head self-attention with trainable projections (Lab 4, Part B).},label={lst:attention:sa}]
class SelfAttention(nn.Module):
    def __init__(self, d_in, d_out, qkv_bias=False):
        super().__init__()
        self.W_query = nn.Linear(d_in, d_out, bias=qkv_bias)
        self.W_key   = nn.Linear(d_in, d_out, bias=qkv_bias)
        self.W_value = nn.Linear(d_in, d_out, bias=qkv_bias)

    def forward(self, x):                       # (..., T, d_in)
        q, k, v = self.W_query(x), self.W_key(x), self.W_value(x)
        scores = q @ k.transpose(-2, -1) / math.sqrt(k.shape[-1])
        return torch.softmax(scores, dim=-1) @ v   # (..., T, d_out)
\end{lstlisting}

A first version might store the projections as raw \code{nn.Parameter(torch.rand(d\_in, d\_out))}
and compute \code{x @ W}. Using \code{nn.Linear} instead changes two things. First, \code{nn.Linear}
stores its weight as $(d_{\text{out}},d_{\text{in}})$ and computes $\vect{x}\mat{W}^\top$---the two
versions compute the same function if one copies \code{W.T} into the linear layer, which is a
useful test. Second, \code{nn.Linear} uses a scaled (Kaiming-uniform) initialisation, whereas
\code{torch.rand} draws from $[0,1)$; all-positive, unscaled weights make scores grow with $d$ and
training less stable.

% =====================================================================
\section{Hiding future words with causal attention}\label{sec:attention:causal}

A language model is trained to predict token $t+1$ from tokens $1..t$. If position $t$ could attend
to position $t+1$ during training, the task would be trivial (copy the answer) and the model would
learn nothing useful for generation, where the future does not yet exist. \emph{Causal} (or
\emph{masked}) attention restricts query $i$ to keys $j\le i$.

\subsection{Applying a causal attention mask}
There are two equivalent ways to apply the mask (\cref{fig:attention:mask}).
\begin{enumerate}
  \item \emph{Mask after softmax, then renormalise}: compute $\mat{A}=\softmax(\mat{S})$, zero the entries
  above the diagonal, and divide each row by its new sum.
  \item \emph{Mask before softmax}: set $s_{ij}=-\infty$ for $j>i$, then apply softmax. Since
  $e^{-\infty}=0$ the masked entries get exactly zero weight and the remaining entries are already
  normalised.
\end{enumerate}
They agree exactly: in route 1, row $i$'s surviving entries are
$e^{s_{ij}}/Z$ with $Z=\sum_{k} e^{s_{ik}}$ over \emph{all} $k$; renormalising divides by
$\sum_{k\le i}e^{s_{ik}}/Z$, and the full-row $Z$ cancels. In particular, the future \emph{does not
leak}: although route 1's first softmax looks at future scores, their influence is removed entirely
by the renormalisation. Route 2 is what every implementation uses---it is one fused
operation, does less work, and is what FlashAttention exploits to skip masked tiles entirely.

\begin{figure}[t]
\centering
\begin{tikzpicture}[x=6mm,y=6mm,
  arr/.style={-{Stealth},thick}]
\newcommand{\grid}[2]{%
  \foreach \r in {0,...,5} \foreach \c in {0,...,5} {
    \pgfmathsetmacro{\ok}{\c<=\r ? 1 : 0}
    \ifnum\c>\r \fill[#2] (#1+\c,-\r) rectangle ++(1,-1); \else \fill[fcompute!25] (#1+\c,-\r) rectangle ++(1,-1); \fi
  }
  \draw[srcgray,step=1] (#1,0) grid (#1+6,-6);}
\grid{0}{fcompute!25}
\node[font=\scriptsize,align=center] at (3,-6.8) {scores $\mat{S}$};
\grid{9}{warnred!35}
\foreach \r in {0,...,4} \foreach \c in {\r,...,4} { \node[font=\tiny] at (9+\c+1.5,-\r-0.5) {$-\infty$}; }
\node[font=\scriptsize,align=center] at (12,-6.8) {masked: $s_{ij}{=}-\infty$ for $j{>}i$};
\grid{18}{white}
\node[font=\scriptsize,align=center] at (21,-6.8) {$\softmax$ rows: zeros above\\diagonal, rows sum to 1};
\draw[arr] (6.3,-3) -- (8.7,-3) node[midway,above,font=\scriptsize] {mask};
\draw[arr] (15.3,-3) -- (17.7,-3) node[midway,above,font=\scriptsize] {softmax};
\end{tikzpicture}
\caption{Causal masking applied before the softmax. Only the lower triangle (including the
diagonal) receives probability mass.}
\label{fig:attention:mask}
\end{figure}

\begin{lstlisting}[caption={Two routes to causal weights; they are equal to float precision.},label={lst:attention:causalroutes}]
def causal_softmax(scores):                     # route 2: -inf before softmax
    T = scores.shape[-1]
    mask = torch.triu(torch.ones(T, T, dtype=torch.bool), diagonal=1)
    return torch.softmax(scores.masked_fill(mask, float("-inf")), dim=-1)

def mask_after_softmax(scores):                 # route 1: zero, renormalise
    w = torch.tril(torch.softmax(scores, dim=-1))
    return w / w.sum(dim=-1, keepdim=True)
\end{lstlisting}

\subsection{Masking additional attention weights with dropout}
\emph{Dropout} zeroes a random subset of activations during training (each with probability $p$)
and scales the survivors by $1/(1-p)$ so expected values are unchanged; at evaluation time it is
the identity. In GPT-2, dropout is applied in three places, one of them to the attention weight
matrix after the softmax, so that a token cannot rely too heavily on any single other position
(\cref{fig:attention:dropout}). With $p=0.5$ a weight of 0.28 becomes either 0 or 0.56, and a row
no longer sums exactly to 1---only in expectation.

\begin{figure}[t]
\centering
\begin{tikzpicture}[x=6.5mm,y=6.5mm]
\foreach \r/\vals in {0/{1.0},1/{.33,.67},2/{.24,.34,.42},3/{.25,.19,.28,.28},4/{.17,.23,.21,.14,.25},5/{.14,.20,.22,.15,.12,.17}} {
  \foreach \v [count=\c from 0] in \vals { \attnheatcell{\c}{\r}{\v}{140} } }
\draw[srcgray] (0,0) rectangle (6,-6);
\node[font=\scriptsize] at (3,-6.6) {causal weights};
\draw[-{Stealth},thick] (6.4,-3) -- (8.6,-3) node[midway,above,font=\scriptsize] {dropout};
\node[font=\scriptsize] at (7.5,-3.6) {$p=0.5$};
\begin{scope}[xshift=58.5mm]
\foreach \r/\vals in {0/{2.0},1/{0,1.34},2/{.48,0,.84},3/{0,.38,0,.56},4/{.34,0,.42,.28,0},5/{0,.40,.44,0,.24,0}} {
  \foreach \v [count=\c from 0] in \vals { \attnheatcell{\c}{\r}{\v}{70} } }
\draw[srcgray] (0,0) rectangle (6,-6);
\node[font=\scriptsize] at (3,-6.6) {survivors scaled by $1/(1-p)=2$};
\end{scope}
\end{tikzpicture}
\caption{Dropout on the attention weights (one random draw, illustrative). Zeros are dropped
connections; surviving weights double so each row's \emph{expected} sum is still 1.}
\label{fig:attention:dropout}
\end{figure}

\begin{systemslens}[Dropout is not free in fused kernels]
In eager PyTorch, attention dropout materialises a $B\times H\times T\times T$ random mask and
stores it for the backward pass---as large as the attention matrix itself. FlashAttention instead
generates dropout bits on the fly from a counter-based RNG (Philox) keyed by a seed and an offset,
and \emph{regenerates} identical bits in the backward pass, storing nothing. This is why fused
kernels take a seed/offset rather than a mask, and why bitwise reproducibility of dropout across
different kernels or tile sizes is not guaranteed. Most modern LLMs pretrain with attention dropout
set to 0 anyway; with trillions of tokens and a single epoch, overfitting is not the concern and
the extra RNG work is wasted.
\end{systemslens}

\subsection{Implementing a compact causal attention class}
Two practical changes turn the above into a reusable layer: it must accept batches of shape
$(B,T,d_{\text{in}})$, and the mask should be created once and moved with the module. Registering the
mask with \code{register\_buffer} does the latter: buffers are moved by \code{.to(device)} with the
parameters but are not trained. We register it with \code{persistent=False} so it is not written to
checkpoints (it can always be recomputed).

\begin{lstlisting}[caption={Batched causal attention with dropout (Lab 4, Part C).},label={lst:attention:causal}]
class CausalAttention(nn.Module):
    def __init__(self, d_in, d_out, context_length, dropout=0.0,
                 qkv_bias=False):
        super().__init__()
        self.d_out = d_out
        self.W_query = nn.Linear(d_in, d_out, bias=qkv_bias)
        self.W_key   = nn.Linear(d_in, d_out, bias=qkv_bias)
        self.W_value = nn.Linear(d_in, d_out, bias=qkv_bias)
        self.dropout = nn.Dropout(dropout)
        mask = torch.triu(torch.ones(context_length, context_length), diagonal=1)
        self.register_buffer("mask", mask.bool(), persistent=False)

    def forward(self, x):                                   # (B, T, d_in)
        B, T, _ = x.shape
        q, k, v = self.W_query(x), self.W_key(x), self.W_value(x)
        scores = q @ k.transpose(1, 2) / math.sqrt(self.d_out)  # (B, T, T)
        scores = scores.masked_fill(self.mask[:T, :T], float("-inf"))
        w = self.dropout(torch.softmax(scores, dim=-1))
        return w @ v                                         # (B, T, d_out)
\end{lstlisting}

The slice \code{self.mask[:T, :T]} lets the layer handle any length up to \code{context\_length}.
A good test of causality (Lab 4) perturbs the inputs at positions $\ge t$ and checks that outputs at
positions $<t$ are bit-for-bit unchanged.

% =====================================================================
\section{Extending single-head attention to multi-head attention}\label{sec:attention:mha}

A single head computes one attention pattern per query. But a token typically needs several
kinds of information at once---the subject it agrees with, the noun a pronoun refers to, the
previous token, the matching bracket. \emph{Multi-head attention} runs $H$ attention heads in
parallel, each with its own projections into a smaller $\dhead=d_{\text{out}}/H$-dimensional
subspace, concatenates their outputs, and mixes them with an output projection $\mat{W}_o$:
\begin{equation}
\operatorname{MHA}(\mat{X}) = \big[\,\operatorname{head}_1\,\Vert\,\cdots\,\Vert\,\operatorname{head}_H\big]\mat{W}_o,\qquad
\operatorname{head}_h = \operatorname{Attention}(\mat{X}\mat{W}_q^{(h)},\mat{X}\mat{W}_k^{(h)},\mat{X}\mat{W}_v^{(h)}).
\end{equation}
The total parameter count and FLOPs match one big head of width $d_{\text{out}}$; what changes is
that the model gets $H$ independent attention distributions per token.

\subsection{Stacking multiple single-head attention layers}
The literal implementation keeps a list of \code{CausalAttention} modules and concatenates their
outputs along the feature dimension. With $d_{\text{out}}=2$ per head and $H=2$ heads, the output is
$(B,T,4)$.

\begin{lstlisting}[caption={Multi-head attention as a list of heads.},label={lst:attention:wrapper}]
class MultiHeadAttentionWrapper(nn.Module):
    def __init__(self, d_in, d_out, context_length, dropout, num_heads,
                 qkv_bias=False):
        super().__init__()
        self.heads = nn.ModuleList(
            [CausalAttention(d_in, d_out, context_length, dropout, qkv_bias)
             for _ in range(num_heads)])

    def forward(self, x):
        return torch.cat([h(x) for h in self.heads], dim=-1)   # (B, T, H*d_out)
\end{lstlisting}

\subsection{Implementing multi-head attention with weight splits}
The wrapper runs the heads one after another: $3H$ small projection GEMMs, $H$ score GEMMs, $H$
softmaxes and $H$ value GEMMs, each a separate kernel launch on a GPU. The standard implementation
instead uses \emph{one} $d_{\text{in}}\times d_{\text{out}}$ projection each for $\mat{Q},\mat{K},\mat{V}$ and
splits the result into heads by reshaping (\cref{fig:attention:split}):
$(B,T,d_{\text{out}})\xrightarrow{\text{view}}(B,T,H,\dhead)\xrightarrow{\text{transpose}}(B,H,T,\dhead)$.
Batched matrix multiplication then treats $(B,H)$ as batch dimensions, so all heads are computed
by single kernels.

\begin{figure}[t]
\centering
\begin{tikzpicture}[
  slab/.style={draw,minimum height=12mm,align=center,font=\scriptsize,inner sep=1pt},
  arr/.style={-{Stealth},thick}]
\node[slab,minimum width=24mm,fill=warnredbg,draw=warnred] (q) {$\mat{Q}$\\$(B,T,d_{\text{out}})$\\\tiny$(2,6,4)$};
\node[slab,minimum width=24mm,right=12mm of q,fill=warnredbg,draw=warnred] (qv) {};
\draw[warnred] ($(qv.north west)!0.5!(qv.north east)$) -- ($(qv.south west)!0.5!(qv.south east)$);
\node[font=\scriptsize] at (qv) {\ \ head 1 \quad head 2};
\node[font=\scriptsize,below=0mm of qv] {$(B,T,H,\dhead)$ \tiny$(2,6,2,2)$};
\draw[arr] (q)--(qv) node[midway,above,font=\scriptsize] {\code{view}};
\node[slab,minimum width=12mm,right=14mm of qv,yshift=8mm,fill=warnred!15,draw=warnred] (h1) {head 1\\$(T,\dhead)$};
\node[slab,minimum width=12mm,right=14mm of qv,yshift=-8mm,fill=warnred!15,draw=warnred] (h2) {head 2\\$(T,\dhead)$};
\draw[arr] (qv.east) -- (h1.west); \draw[arr] (qv.east) -- (h2.west);
\node[font=\scriptsize,above=0mm of h1] {\code{transpose(1,2)}: $(B,H,T,\dhead)$};
\node[slab,minimum width=20mm,right=12mm of h1,fill=labpurplebg,draw=fcomm] (a1) {$\softmax(\frac{\mat{Q}_1\mat{K}_1^\top}{\sqrt{\dhead}})\mat{V}_1$};
\node[slab,minimum width=20mm,right=12mm of h2,fill=labpurplebg,draw=fcomm] (a2) {$\softmax(\frac{\mat{Q}_2\mat{K}_2^\top}{\sqrt{\dhead}})\mat{V}_2$};
\draw[arr] (h1)--(a1); \draw[arr] (h2)--(a2);
\node[slab,minimum width=22mm,below=14mm of qv,fill=napkinbg,draw=fmemory] (cat) {merge heads\\$(B,T,d_{\text{out}})$};
\node[slab,minimum width=22mm,left=12mm of cat,fill=keygreenbg,draw=fdata] (out) {$\cdot\,\mat{W}_o$\\$(B,T,d_{\text{out}})$};
\draw[arr] (a2.south) |- (cat.east); \draw[arr] (cat)--(out);
\node[font=\scriptsize,text=srcgray,below=0mm of cat] {\code{transpose(1,2).contiguous().view(B,T,d)}};
\end{tikzpicture}
\caption{Weight-split multi-head attention: one projection, then reshape into heads, batched
attention over $(B,H)$, merge, output projection. Same shapes for $\mat{K}$ and $\mat{V}$.}
\label{fig:attention:split}
\end{figure}

\begin{lstlisting}[caption={Multi-head attention with weight splits (Lab 4, Part D; same parameter names as \texttt{llmsys.layers.MultiHeadAttention}).},label={lst:attention:mha}]
class MultiHeadAttention(nn.Module):
    def __init__(self, d_in, d_out, context_length, dropout, num_heads,
                 qkv_bias=False):
        super().__init__()
        assert d_out % num_heads == 0
        self.d_out, self.num_heads = d_out, num_heads
        self.head_dim = d_out // num_heads
        self.W_query = nn.Linear(d_in, d_out, bias=qkv_bias)
        self.W_key   = nn.Linear(d_in, d_out, bias=qkv_bias)
        self.W_value = nn.Linear(d_in, d_out, bias=qkv_bias)
        self.out_proj = nn.Linear(d_out, d_out)      # mixes the heads
        self.dropout = nn.Dropout(dropout)
        mask = torch.triu(torch.ones(context_length, context_length), diagonal=1)
        self.register_buffer("mask", mask.bool(), persistent=False)

    def forward(self, x):
        B, T, _ = x.shape
        H, Dh = self.num_heads, self.head_dim
        split = lambda t: t.view(B, T, H, Dh).transpose(1, 2)    # (B, H, T, Dh)
        q, k, v = (split(f(x)) for f in (self.W_query, self.W_key, self.W_value))
        scores = q @ k.transpose(2, 3) / math.sqrt(Dh)          # (B, H, T, T)
        scores = scores.masked_fill(self.mask[:T, :T], float("-inf"))
        w = self.dropout(torch.softmax(scores, dim=-1))
        ctx = (w @ v).transpose(1, 2).contiguous().view(B, T, self.d_out)
        return self.out_proj(ctx)
\end{lstlisting}

Note the \code{.contiguous()} before the final \code{view}: after \code{transpose} the tensor's
memory layout no longer matches its logical shape, and \code{view} requires compatible strides.
\code{.contiguous()} performs a physical copy into the merged-head layout (alternatively,
\code{.reshape} copies only when needed). For GPT-2 small, $H=12$ heads of $\dhead=64$ give
$d_{\text{out}}=768$; GPT-2 XL uses 25 heads of 64 for $d=1600$.

\begin{systemslens}[Why weight splitting is faster: kernel granularity]
For GPT-2 small ($d=768$, $H=12$) the wrapper performs 36 projections of size
$(BT\times768)\cdot(768\times64)$; the weight-split version performs 3 of size
$(BT\times768)\cdot(768\times768)$---the same FLOPs. But a $768\times64$ weight gives each GEMM far
fewer output tiles to spread over a GPU's 100+ streaming multiprocessors, so many SMs idle, and each
call pays a few microseconds of launch overhead. GPT-2's original code goes one step further and
fuses Q, K and V into a single $768\times2304$ projection (\code{c\_attn}), which is why the weight
loader in \code{llmsys/gpt.py} splits it in three. The general lesson---\emph{batch small
operations into large ones}---reappears at every level of this book, from heads to requests
(\cref{ch:serving}).
\end{systemslens}

% =====================================================================
\section{The cost of attention}\label{sec:attention:cost}

We now treat attention as a workload. Let $B$ be the batch size, $T$ the sequence length, $H$ the
number of heads, $\dhead$ the head dimension and $d=H\dhead$ the model width.

\paragraph{FLOPs.} Counting a multiply--add as 2 FLOPs, one forward pass of one attention layer costs
\begin{align}
\text{projections } (\mat{Q},\mat{K},\mat{V},\mat{W}_o) &: \; 4\cdot 2BTd^2 = 8BTd^2,\\
\mat{Q}\mat{K}^\top \text{ and } \mat{A}\mat{V} &: \; 2\cdot 2BHT^2\dhead = 4BT^2d .
\end{align}
The quadratic term overtakes the projections when $4BT^2d > 8BTd^2$, i.e.\ when $T>2d$: around
1.5k tokens for GPT-2 small, 8k for a $d=4096$ model. Causal masking halves the useful work of the
quadratic term (only the lower triangle matters), which a good kernel exploits.

\Cref{tab:attention:gpt2} instantiates these formulas for one GPT-2-small attention layer. Even at
GPT-2's modest context of 1024, the score and value products are a third of the layer's FLOPs, and the
materialised attention matrix is the largest activation in the layer by a wide margin.

\begin{table}[t]
\centering\small
\begin{tabular}{@{}lrrr@{}}
\toprule
quantity (one layer, $B=8$, $T=1024$, $d=768$, $H=12$) & formula & value & share \\
\midrule
Q, K, V, O projections (fwd FLOPs) & $8BTd^2$ & 38.7 GFLOP & 67\% \\
$\mat{Q}\mat{K}^\top$ and $\mat{A}\mat{V}$ (fwd FLOPs, dense) & $4BT^2d$ & 25.8 GFLOP & 33\% \\
\midrule
input / output activations (fp32) & $BTd\cdot4$\,B & 25.2 MB each & \\
$\mat{Q},\mat{K},\mat{V}$ (fp32) & $3BTd\cdot4$\,B & 75.5 MB & \\
attention matrix $\mat{A}$ (fp32), kept for backward & $BHT^2\cdot4$\,B & 403 MB & \\
\bottomrule
\end{tabular}
\caption{FLOPs and activation memory of one GPT-2-small attention layer (analytic). The $T^2$ terms
are already significant at $T=1024$; they dominate beyond $T\approx2d$.}
\label{tab:attention:gpt2}
\end{table}

\paragraph{Memory.} The naive implementation materialises the scores $\mat{S}$ and probabilities
$\mat{A}$, each $B\cdot H\cdot T^2$ elements, and keeps $\mat{A}$ (and a dropout mask) for the backward
pass. Everything else is $O(BTd)$.

\begin{napkinmath}[The $T^2$ wall]
Take one layer of a model with $H=32$ heads at $T=8192$, $B=1$, bf16 (2 bytes):
$32\cdot 8192^2\cdot 2 = 4.3$\,GB for one copy of the attention matrix. Training must keep it for
backward in \emph{every} layer: 32 layers would need $\approx 137$\,GB---more than an 80\,GB H100
holds, for a single sequence. At $T=32{,}768$ one layer alone needs $\approx 69$\,GB.
Long-context training with materialised attention is simply impossible.
\end{napkinmath}

\paragraph{Arithmetic intensity.} Is naive attention compute-bound or memory-bound? Per head,
$\mat{Q}\mat{K}^\top$ performs $2T^2\dhead$ FLOPs and must at minimum write $T^2$ scores. The
mask, scaling, softmax and dropout are elementwise passes over $T^2$ values with only a handful of
FLOPs each. With $\dhead=128$ the score GEMM itself is fine, but the elementwise passes have an
arithmetic intensity of order 1 FLOP/byte, far below the hundreds of FLOPs/byte an H100 needs to be
compute-bound (\cref{ch:pytorch}).

\begin{napkinmath}[Naive attention on an H100: where the time goes]
One head, $T=8192$, $\dhead=128$, bf16. Matmul FLOPs: $4T^2\dhead = 34.4$\,GFLOP, i.e.\
$\approx 35\,\mu$s at 989\,TFLOP/s. If the unfused implementation makes about four round trips of
the $T\times T$ matrix through HBM (write $\mat{S}$; read $\mat{S}$, write $\mat{A}$; read $\mat{A}$
for $\mat{A}\mat{V}$), traffic is $4\cdot T^2\cdot 2\,\text{B}\approx 537$\,MB, i.e.\
$\approx 160\,\mu$s at 3.35\,TB/s. The kernel spends about $4.6\times$ longer moving the attention
matrix than computing with it. Avoiding those round trips---not reducing FLOPs---is the goal.
\end{napkinmath}

\begin{figure}[t]
\centering
\begin{tikzpicture}
\begin{axis}[width=0.8\linewidth,height=6cm,xmode=log,ymode=log,log basis x=2,
  xlabel={sequence length $T$},ylabel={bytes per layer (B$=$1, H$=$32, bf16)},
  xmin=512,xmax=131072,ymin=1e5,ymax=1e13,legend pos=south east,legend style={font=\scriptsize},
  font=\small,grid=major,grid style={gray!15}]
\addplot[thick,warnred,domain=512:131072,samples=30] {32*x^2*2};
\addlegendentry{naive: $\mat{A}$ stored ($BHT^2$)}
\addplot[thick,fdata,domain=512:131072,samples=30] {32*x*4*2 + 32*x*128*2*4};
\addlegendentry{tiled: $\mat{Q},\mat{K},\mat{V},\mat{O}$ + log-sum-exp ($O(BHT\dhead)$)}
\addplot[dashed,black,domain=512:131072,samples=2] {80e9};
\node[font=\scriptsize,anchor=south east] at (axis cs:131072,80e9) {80\,GB (one H100)};
\end{axis}
\end{tikzpicture}
\caption{Memory needed per attention layer for the backward pass (analytic). The naive
implementation stores the $T\times T$ probabilities; IO-aware attention stores only inputs, output
and one scalar per row.}
\label{fig:attention:mem}
\end{figure}

% =====================================================================
\section{IO-aware attention: the online softmax and FlashAttention}\label{sec:attention:flash}

FlashAttention (Dao et al., 2022) computes \emph{exact} attention---the same numbers as
\cref{eq:attention:sdpa} up to floating-point reassociation---without writing $\mat{S}$ or
$\mat{A}$ to HBM. It rests on two ideas: \emph{tiling}, so that blocks of $\mat{Q}$, $\mat{K}$, $\mat{V}$
fit in on-chip SRAM (shared memory and registers), and the \emph{online softmax}, which makes it
possible to normalise rows whose entries are seen one block at a time.

\subsection{The online softmax}
The softmax of a row $\vect{x}$ needs its maximum $m=\max_j x_j$ (for stability) and normaliser
$\ell=\sum_j e^{x_j-m}$---global quantities over the whole row. The trick, due to Milakov and
Gimelshein (2018), is that both can be maintained incrementally. Suppose we have processed a prefix
of the row and hold $(m,\ell)$ for it. A new block $\vect{x}_B$ arrives:
\begin{equation}
m' = \max\big(m,\ \max_j x_{B,j}\big),\qquad
\ell' = \ell\,e^{m-m'} + \sum_j e^{x_{B,j}-m'} .
\label{eq:attention:online}
\end{equation}
The factor $e^{m-m'}\le 1$ \emph{rescales} the old normaliser from the old reference maximum to
the new one. After the last block, $(m,\ell)$ are exactly the global values.

\begin{lstlisting}[caption={Online softmax over blocks (Lab 4, Part E).},label={lst:attention:onlinesm}]
def online_softmax(x, block):
    m = torch.full(x.shape[:-1], float("-inf"))     # running max
    l = torch.zeros(x.shape[:-1])                   # running normaliser
    for s in range(0, x.shape[-1], block):
        xb = x[..., s:s + block]
        m_new = torch.maximum(m, xb.max(dim=-1).values)
        l = l * torch.exp(m - m_new) \
            + torch.exp(xb - m_new[..., None]).sum(-1)
        m = m_new
    return torch.exp(x - m[..., None]) / l[..., None]   # second pass
\end{lstlisting}

Listing~\ref{lst:attention:onlinesm} still makes a second pass to emit probabilities. Attention does
not need the probabilities themselves, only $\mat{A}\mat{V}$, and the same rescaling applies to a
running \emph{output accumulator}: if $\vect{o}$ holds $\sum_j e^{s_j-m}\vect{v}_j$ for the prefix,
then on a new block $\vect{o}' = e^{m-m'}\vect{o} + \sum_{j\in B}e^{s_j-m'}\vect{v}_j$, and the final
output is $\vect{o}/\ell$. A single pass over $\mat{K}$ and $\mat{V}$ suffices.

\begin{figure}[t]
\centering
\begin{tikzpicture}[
  blk/.style={draw,minimum width=9mm,minimum height=9mm,font=\scriptsize,inner sep=0},
  arr/.style={-{Stealth},thick}]
% HBM
\node[draw=srcgray,dashed,rounded corners,minimum width=86mm,minimum height=30mm,label={[font=\small\bfseries,text=srcgray]above:HBM (large, slow: 3.35 TB/s)}] (hbm) at (0,0) {};
\node[blk,fill=warnredbg,draw=warnred,minimum height=24mm,minimum width=8mm] (Q) at (-3.4,0) {};
\foreach \i in {0,...,3} \draw[warnred] (-3.8,{1.2-\i*0.6-0.6}) -- (-3.0,{1.2-\i*0.6-0.6});
\node[font=\scriptsize] at (-3.4,-1.45) {$\mat{Q}$};
\node[blk,fill=sysbluebg,draw=fcompute,minimum height=24mm,minimum width=8mm] (K) at (-1.6,0) {};
\foreach \i in {0,...,3} \draw[fcompute] (-2.0,{1.2-\i*0.6-0.6}) -- (-1.2,{1.2-\i*0.6-0.6});
\node[font=\scriptsize] at (-1.6,-1.45) {$\mat{K}$};
\node[blk,fill=keygreenbg,draw=fdata,minimum height=24mm,minimum width=8mm] (V) at (0.2,0) {};
\foreach \i in {0,...,3} \draw[fdata] (-0.2,{1.2-\i*0.6-0.6}) -- (0.6,{1.2-\i*0.6-0.6});
\node[font=\scriptsize] at (0.2,-1.45) {$\mat{V}$};
\node[blk,fill=napkinbg,draw=fmemory,minimum height=24mm,minimum width=8mm] (O) at (2.2,0) {};
\node[font=\scriptsize] at (2.2,-1.45) {$\mat{O}$, $\vect{L}$};
\node[font=\scriptsize,text=srcgray,align=center] at (3.6,0) {no $T{\times}T$\\matrix here};
% SRAM
\node[draw=fcomm,thick,rounded corners,minimum width=86mm,minimum height=26mm,label={[font=\small\bfseries,text=fcomm]below:on-chip SRAM per SM (tiny, fast)}] (sram) at (0,-4.4) {};
\node[blk,fill=warnredbg,draw=warnred,minimum height=6mm] (qt) at (-3.4,-4.0) {$\mat{Q}_i$};
\node[blk,fill=sysbluebg,draw=fcompute,minimum height=6mm] (kt) at (-2.2,-4.0) {$\mat{K}_j$};
\node[blk,fill=keygreenbg,draw=fdata,minimum height=6mm] (vt) at (-1.0,-4.0) {$\mat{V}_j$};
\node[blk,fill=labpurplebg,draw=fcomm,minimum width=14mm,minimum height=6mm] (st) at (0.7,-4.0) {$\mat{S}_{ij},\mat{P}_{ij}$};
\node[blk,fill=napkinbg,draw=fmemory,minimum width=17mm,minimum height=6mm] (acc) at (2.75,-4.0) {$\vect{m}_i,\vect{\ell}_i,\mat{O}_i$ acc.};
\draw[arr,warnred] (-3.4,-1.2) -- (qt) node[midway,left,font=\tiny] {once};
\draw[arr,fcompute] (-1.6,-1.2) -- (kt) node[midway,left,font=\tiny] {stream $j$};
\draw[arr,fdata] (0.2,-1.2) -- (vt);
\draw[arr,fmemory] (acc.north) -- (2.2,-1.2) node[midway,right,font=\tiny] {once};
\node[font=\scriptsize,text=srcgray,align=left,anchor=west] at (-4.2,-4.95) {for each key/value tile $j$: $\mat{S}=\mat{Q}_i\mat{K}_j^\top$, update $(\vect{m},\vect{\ell},\mat{O})$ by \eqref{eq:attention:online}};
\end{tikzpicture}
\caption{FlashAttention's data movement for one query tile $\mat{Q}_i$ (one thread block). $\mat{Q}_i$
is loaded once, key/value tiles stream through SRAM, scores and probabilities live only on chip,
and the output tile plus one log-sum-exp per row are written once.}
\label{fig:attention:tiling}
\end{figure}

\subsection{The tiled forward pass}
Putting the pieces together gives Algorithm~4.1 (boxed below), implemented in plain PyTorch in
\code{llmsys/flash.py} (\cref{lst:attention:flashfwd}). On a GPU, the outer loop over query tiles is
the grid of thread blocks (one block per $(b,h,i)$), and the inner loop runs inside each block,
keeping $\mat{Q}_i$, the current $\mat{K}_j,\mat{V}_j$, the score tile and the accumulators in
shared memory and registers. Appendix~\ref{app:triton} writes this kernel in Triton.

\begin{lstlisting}[caption={Tiled attention forward with online softmax (\texttt{llmsys/flash.py}, simplified).},label={lst:attention:flashfwd}]
for qs in range(0, Tq, block_q):                       # parallel on a GPU
    qb = q[..., qs:qs + block_q, :]
    m = full(-inf); l = zeros(); acc = zeros(block_q, D)
    k_end = min(Tk, qs + block_q + offset) if causal else Tk  # skip future
    for ks in range(0, k_end, block_k):                # sequential per block
        kb, vb = k[..., ks:ks + block_k, :], v[..., ks:ks + block_k, :]
        s = (qb @ kb.transpose(-2, -1)) * scale        # lives in SRAM only
        if causal:
            s = s.masked_fill(k_pos[None, :] > q_pos[:, None], -inf)
        m_new = torch.maximum(m, s.max(-1).values)
        p = torch.exp(s - m_new[..., None])
        alpha = torch.exp(m - m_new)                   # rescale old state
        l = alpha * l + p.sum(-1)
        acc = alpha[..., None] * acc + p @ vb
        m = m_new
    out[..., qs:qs + block_q, :] = acc / l[..., None]
    lse[..., qs:qs + block_q] = m + torch.log(l)       # saved for backward
\end{lstlisting}

\begin{figure}[t]
\centering\small
\fbox{\begin{minipage}{0.92\linewidth}
\textbf{Algorithm 4.1} (tiled attention forward; one $(b,h)$ slice)\\[2pt]
\textbf{Input:} $\mat{Q}\in\R^{T_q\times D}$, $\mat{K},\mat{V}\in\R^{T_k\times D}$, tile sizes $B_r,B_c$.\\
\textbf{for} each query tile $\mat{Q}_i$ ($B_r$ rows) \textbf{in parallel}:\\
\hspace*{1.2em} $\vect{m}\gets-\infty$, $\vect{\ell}\gets 0$, $\mat{O}\gets 0$ \hfill (registers/SRAM)\\
\hspace*{1.2em} \textbf{for} each key/value tile $j$ not entirely in the causal future:\\
\hspace*{2.4em} load $\mat{K}_j,\mat{V}_j$; \ $\mat{S}\gets\mat{Q}_i\mat{K}_j^\top/\sqrt{D}$ (+ mask)\\
\hspace*{2.4em} $\vect{m}'\gets\max(\vect{m},\operatorname{rowmax}\mat{S})$; \ $\mat{P}\gets e^{\mat{S}-\vect{m}'}$; \ $\vect{\alpha}\gets e^{\vect{m}-\vect{m}'}$\\
\hspace*{2.4em} $\vect{\ell}\gets\vect{\alpha}\odot\vect{\ell}+\operatorname{rowsum}\mat{P}$; \ $\mat{O}\gets\operatorname{diag}(\vect{\alpha})\mat{O}+\mat{P}\mat{V}_j$; \ $\vect{m}\gets\vect{m}'$\\
\hspace*{1.2em} write $\mat{O}_i\gets\operatorname{diag}(\vect{\ell})^{-1}\mat{O}$ and $\vect{L}_i\gets\vect{m}+\log\vect{\ell}$ to HBM.
\end{minipage}}
\end{figure}

\paragraph{IO complexity.} Let $M$ be the SRAM capacity in elements. A query tile of $B_r$ rows needs
room for itself, a key and a value tile, and the accumulator, so $B_r,B_c=\Theta(M/D)$. Each of the
$T/B_r$ query tiles streams all of $\mat{K}$ and $\mat{V}$ once, for $\Theta(T^2D/B_r)=\Theta(T^2D^2/M)$
element reads, versus $\Theta(TD+T^2)$ for the naive algorithm. Dao et al.\ prove this is optimal up
to constants for exact attention. Because $D^2/M$ is a small constant for typical head dimensions
(about $1/6$ for $D=128$ in bf16 with H100-sized shared memory), the asymptotic win is a constant
factor of traffic \emph{plus} the elimination of all $T^2$-sized intermediate tensors and
the separate mask/softmax/dropout kernels; in practice the latter dominate.

\subsection{The backward pass: recompute instead of store}
Backpropagation through $\mat{O}=\mat{P}\mat{V}$ with $\mat{P}=\softmax(\mat{S})$ normally needs $\mat{P}$.
FlashAttention stores only $\mat{O}$ and the per-row log-sum-exp $\vect{L}$, and \emph{recomputes}
each tile $\mat{P}_{ij}=\exp(\mat{S}_{ij}-\vect{L}_i)$ during the backward pass from $\mat{Q}_i$ and
$\mat{K}_j$. With $\mat{dO}$ the incoming gradient, the tile-local gradients are
\begin{gather}
\mat{dV}_j \mathrel{+}= \mat{P}_{ij}^\top\mat{dO}_i,\qquad
\mat{dP}_{ij}=\mat{dO}_i\mat{V}_j^\top,\qquad
\mat{dS}_{ij}=\mat{P}_{ij}\odot(\mat{dP}_{ij}-\vect{\delta}_i),\\
\mat{dQ}_i \mathrel{+}= \mat{dS}_{ij}\mat{K}_j,\qquad
\mat{dK}_j \mathrel{+}= \mat{dS}_{ij}^\top\mat{Q}_i ,
\end{gather}
(each $\mat{dS}$ also carries the $1/\sqrt{D}$ scale), where the softmax Jacobian's row term
$\vect{\delta}_i=\operatorname{rowsum}(\mat{P}_i\odot\mat{dP}_i)$ would seem to need the whole row. It
does not: $\operatorname{rowsum}(\mat{P}\odot\mat{dO}\mat{V}^\top)=\operatorname{rowsum}(\mat{dO}\odot\mat{P}\mat{V})=\operatorname{rowsum}(\mat{dO}\odot\mat{O})$,
which is computed once up front from $\mat{O}$ and $\mat{dO}$. Recomputation costs roughly one extra
$\mat{Q}\mat{K}^\top$ GEMM, but saves storing and re-reading $O(T^2)$ bytes; on memory-bound hardware this
makes the backward pass \emph{faster}, not just leaner. This is activation checkpointing
(\cref{ch:pretrain}) applied at the granularity of a single kernel.

\begin{systemslens}[What FlashAttention-2 and -3 changed]
FlashAttention-2 (Dao, 2023) reorganised the work: it parallelises over the sequence dimension as
well as batch and heads (important when $B\cdot H$ is smaller than the number of SMs, as in
long-context training), swaps the loop order in the backward pass so that $\mat{dK}_j,\mat{dV}_j$ are
accumulated in registers and written once, partitions work across warps to avoid shared-memory
communication, and minimises non-matmul FLOPs (e.g.\ deferring the $1/\ell$ rescale to the end),
because on tensor-core GPUs a non-matmul FLOP is an order of magnitude more expensive than a matmul
FLOP. FlashAttention-3 (Shah et al., 2024) targets Hopper: asynchronous TMA loads and
warp-specialised producer/consumer pipelines overlap data movement with WGMMA tensor-core
instructions, softmax of one tile is overlapped with the GEMMs of the next, and FP8 support uses
incoherent processing to limit quantisation error. For decoding, where $T_q=1$, parallelising over
queries is useless; \emph{split-KV} kernels (FlashDecoding) parallelise over chunks of the KV cache
and combine partial results with the same log-sum-exp algebra (\cref{ch:inference}).
\end{systemslens}

\begin{pitfall}[Exact, but not bit-identical]
Tiling changes the order of floating-point additions, so FlashAttention's output differs from the
naive result in the last bits, and different tile sizes or GPUs give different bits. Our float64
tests in Lab 4 agree to $\sim10^{-15}$; in bf16 expect relative differences around $10^{-3}$. This
matters for RL (\cref{ch:rlsys}), where the rollout engine and the trainer may use different
attention kernels and so assign slightly different probabilities to the same tokens.
\end{pitfall}

\subsection{Using it from PyTorch}
In practice one calls PyTorch's fused entry point,
\begin{lstlisting}[numbers=none]
F.scaled_dot_product_attention(q, k, v, is_causal=True)   # "SDPA"
\end{lstlisting}
\noindent and lets PyTorch dispatch to a FlashAttention, memory-efficient
(Rabe and Staats, 2021, style), cuDNN, or plain math backend depending on hardware, dtype, head
dimension and masking. The core \code{llmsys} attention layers accept \code{use\_sdpa=True} to switch
from the explicit implementation to SDPA, and the core test suite checks that both produce the same
logits. One caveat: an arbitrary boolean \code{attn\_mask} can force a slower backend than
\code{is\_causal=True}, because fused kernels generate structured masks on the fly instead of reading
a $T\times T$ mask from memory.


% =====================================================================
\section{Attention at inference time: a preview of the KV cache}\label{sec:attention:decode}

During training, all $T$ queries are processed at once. During generation (\cref{ch:inference}),
the model produces one token per step, and causal masking guarantees that the keys and values of
earlier positions never change once computed. They can therefore be stored---the \emph{KV
cache}---so that each new step projects only the newest token into $\vect{q},\vect{k},\vect{v}$,
appends $\vect{k},\vect{v}$ to the cache, and attends with a single query over all cached keys
(\cref{fig:attention:decode}). The core \code{llmsys} attention modules accept a cache object for
exactly this purpose, and the general causal mask \code{llmsys.layers.causal\_mask(q\_len, k\_len)}
handles $T_q<T_k$ by placing query $i$ at absolute position $T_k-T_q+i$.

\begin{figure}[t]
\centering
\begin{tikzpicture}[
  kv/.style={draw=fcompute,fill=sysbluebg,minimum width=6mm,minimum height=6mm,font=\tiny,inner sep=0},
  nw/.style={kv,draw=warnred,fill=warnredbg},
  arr/.style={-{Stealth},thick}]
\node[font=\small\bfseries,anchor=west] at (-0.4,0.9) {prefill (prompt of 5 tokens)};
\foreach \i in {1,...,5} \node[kv] at ({\i*0.7},0) {$k_\i$};
\foreach \i in {1,...,5} \node[kv,fill=keygreenbg,draw=fdata] at ({\i*0.7},-0.7) {$v_\i$};
\node[font=\scriptsize,align=left,anchor=west] at (4.1,-0.35) {$T_q=T_k=5$: a GEMM,\\compute-bound};
\begin{scope}[yshift=-2.4cm]
\node[font=\small\bfseries,anchor=west] at (-0.4,0.9) {decode step 3 (generating token 8)};
\foreach \i in {1,...,7} \node[kv] at ({\i*0.7},0) {$k_\i$};
\node[nw] at (8*0.7,0) {$k_8$};
\foreach \i in {1,...,7} \node[kv,fill=keygreenbg,draw=fdata] at ({\i*0.7},-0.7) {$v_\i$};
\node[nw] at (8*0.7,-0.7) {$v_8$};
\node[draw=warnred,fill=warnredbg,minimum size=6mm,font=\tiny] (q) at (8*0.7,1.0) {$q_8$};
\foreach \i in {1,...,8} \draw[-{Stealth},fcomm!60] (q.south) -- ({\i*0.7},0.3);
\node[font=\scriptsize,align=left,anchor=west] at (6.4,-0.35) {$T_q=1$, $T_k=8$: a GEMV over the cache,\\memory-bound; only $k_8,v_8$ are new};
\end{scope}
\end{tikzpicture}
\caption{Prefill processes the prompt in parallel and fills the cache; each decode step adds one
key/value pair (red) and attends with a single query over everything cached.}
\label{fig:attention:decode}
\end{figure}

\begin{napkinmath}[Decode attention is a memory stream]
For one decode step with context $T_k$, one layer reads the whole cache: $2\cdot H_{kv}\cdot T_k\cdot\dhead$
values, and does $\approx 4\cdot H\cdot T_k\cdot\dhead$ FLOPs with them. With $H=H_{kv}$ (plain MHA) and
bf16 that is about $4\,\text{FLOPs}/(4\,\text{bytes}) = 1$ FLOP/byte, two to three orders of magnitude
below the H100's balance point of $989/3.35\approx 295$ FLOP/byte. Decode attention runs at memory
bandwidth, so its time is (cache bytes)/(3.35\,TB/s). Two levers follow: read fewer bytes (GQA/MQA,
cache quantisation; \cref{ch:modern,ch:efficient}), or amortise each byte over more queries---which
is what batching many requests does \emph{not} achieve for attention, since each request has its own
cache (\cref{ch:serving}).
\end{napkinmath}

\begin{deepdive}[What do heads learn?]
Interpretability work has found recurring head types in trained transformers: \emph{previous-token
heads} that attend to position $i-1$; \emph{induction heads} (Olsson et al., 2022) that look for an
earlier occurrence of the current token and attend to the token that followed it, implementing
in-context copying of patterns (``\dots A B \dots A $\to$ B''); and heads whose attention collapses
onto the first token or onto delimiters, acting as a no-op ``sink''. Such \emph{attention sinks}
matter for systems: StreamingLLM (Xiao et al., 2024) shows that a windowed KV cache must keep the
first few tokens to remain stable, a constraint any cache-eviction policy must respect.
\end{deepdive}

% =====================================================================
\section{Beyond dense multi-head attention}\label{sec:attention:variants}

The attention of this chapter is \emph{dense}: every query sees every allowed key, and every head has
its own keys and values. Most of the attention variants in modern models change one of those two
properties, for systems reasons more than modelling reasons.

\paragraph{Sharing keys and values: MQA and GQA.} At inference time the keys and values of all past
tokens are cached (\cref{ch:inference}), and reading that cache is what makes decoding slow.
Multi-query attention (Shazeer, 2019) keeps $H$ query heads but a single key/value head; grouped-query
attention (Ainslie et al., 2023) uses $G$ key/value heads, each shared by $H/G$ query heads. The
KV cache shrinks by $H/G$. Qwen3-0.6B has 16 query heads and 8 KV heads; \cref{ch:modern} implements
GQA and quantifies the savings. DeepSeek's multi-head latent attention (MLA) instead caches a
low-rank compressed latent from which keys and values are reconstructed.

\paragraph{Restricting the pattern.} Sliding-window attention lets each token attend only to the
previous $w$ tokens, making cost $O(Tw)$ and the KV cache bounded; models often interleave windowed
and global layers. Block-sparse and dilated patterns (Child et al., 2019; Beltagy et al., 2020) are
other points in the design space. Document masks for packed training sequences
(\cref{ch:text}) are a block-diagonal restriction whose purpose is correctness rather than speed.
Prefix-LM masks let a prompt attend bidirectionally while the continuation remains causal.
\Cref{fig:attention:patterns} shows these masks side by side; all are supported efficiently by
FlashAttention-style kernels because whole tiles can be skipped.

\begin{figure}[t]
\centering
\begin{tikzpicture}[x=3.2mm,y=3.2mm]
\begin{scope}[xshift=0mm]
\fill[fcompute!55] (0,-0) rectangle ++(1,-1) (0,-1) rectangle ++(1,-1) (1,-1) rectangle ++(1,-1) (0,-2) rectangle ++(1,-1) (1,-2) rectangle ++(1,-1) (2,-2) rectangle ++(1,-1) (0,-3) rectangle ++(1,-1) (1,-3) rectangle ++(1,-1) (2,-3) rectangle ++(1,-1) (3,-3) rectangle ++(1,-1) (0,-4) rectangle ++(1,-1) (1,-4) rectangle ++(1,-1) (2,-4) rectangle ++(1,-1) (3,-4) rectangle ++(1,-1) (4,-4) rectangle ++(1,-1) (0,-5) rectangle ++(1,-1) (1,-5) rectangle ++(1,-1) (2,-5) rectangle ++(1,-1) (3,-5) rectangle ++(1,-1) (4,-5) rectangle ++(1,-1) (5,-5) rectangle ++(1,-1) (0,-6) rectangle ++(1,-1) (1,-6) rectangle ++(1,-1) (2,-6) rectangle ++(1,-1) (3,-6) rectangle ++(1,-1) (4,-6) rectangle ++(1,-1) (5,-6) rectangle ++(1,-1) (6,-6) rectangle ++(1,-1) (0,-7) rectangle ++(1,-1) (1,-7) rectangle ++(1,-1) (2,-7) rectangle ++(1,-1) (3,-7) rectangle ++(1,-1) (4,-7) rectangle ++(1,-1) (5,-7) rectangle ++(1,-1) (6,-7) rectangle ++(1,-1) (7,-7) rectangle ++(1,-1) (0,-8) rectangle ++(1,-1) (1,-8) rectangle ++(1,-1) (2,-8) rectangle ++(1,-1) (3,-8) rectangle ++(1,-1) (4,-8) rectangle ++(1,-1) (5,-8) rectangle ++(1,-1) (6,-8) rectangle ++(1,-1) (7,-8) rectangle ++(1,-1) (8,-8) rectangle ++(1,-1) (0,-9) rectangle ++(1,-1) (1,-9) rectangle ++(1,-1) (2,-9) rectangle ++(1,-1) (3,-9) rectangle ++(1,-1) (4,-9) rectangle ++(1,-1) (5,-9) rectangle ++(1,-1) (6,-9) rectangle ++(1,-1) (7,-9) rectangle ++(1,-1) (8,-9) rectangle ++(1,-1) (9,-9) rectangle ++(1,-1) (0,-10) rectangle ++(1,-1) (1,-10) rectangle ++(1,-1) (2,-10) rectangle ++(1,-1) (3,-10) rectangle ++(1,-1) (4,-10) rectangle ++(1,-1) (5,-10) rectangle ++(1,-1) (6,-10) rectangle ++(1,-1) (7,-10) rectangle ++(1,-1) (8,-10) rectangle ++(1,-1) (9,-10) rectangle ++(1,-1) (10,-10) rectangle ++(1,-1) (0,-11) rectangle ++(1,-1) (1,-11) rectangle ++(1,-1) (2,-11) rectangle ++(1,-1) (3,-11) rectangle ++(1,-1) (4,-11) rectangle ++(1,-1) (5,-11) rectangle ++(1,-1) (6,-11) rectangle ++(1,-1) (7,-11) rectangle ++(1,-1) (8,-11) rectangle ++(1,-1) (9,-11) rectangle ++(1,-1) (10,-11) rectangle ++(1,-1) (11,-11) rectangle ++(1,-1);
\draw[srcgray] (0,0) rectangle (12,-12);
\node[font=\scriptsize,align=center,anchor=north] at (6,-12.4) {causal};
\end{scope}
\begin{scope}[xshift=40mm]
\fill[fcompute!55] (0,-0) rectangle ++(1,-1) (0,-1) rectangle ++(1,-1) (1,-1) rectangle ++(1,-1) (0,-2) rectangle ++(1,-1) (1,-2) rectangle ++(1,-1) (2,-2) rectangle ++(1,-1) (0,-3) rectangle ++(1,-1) (1,-3) rectangle ++(1,-1) (2,-3) rectangle ++(1,-1) (3,-3) rectangle ++(1,-1) (1,-4) rectangle ++(1,-1) (2,-4) rectangle ++(1,-1) (3,-4) rectangle ++(1,-1) (4,-4) rectangle ++(1,-1) (2,-5) rectangle ++(1,-1) (3,-5) rectangle ++(1,-1) (4,-5) rectangle ++(1,-1) (5,-5) rectangle ++(1,-1) (3,-6) rectangle ++(1,-1) (4,-6) rectangle ++(1,-1) (5,-6) rectangle ++(1,-1) (6,-6) rectangle ++(1,-1) (4,-7) rectangle ++(1,-1) (5,-7) rectangle ++(1,-1) (6,-7) rectangle ++(1,-1) (7,-7) rectangle ++(1,-1) (5,-8) rectangle ++(1,-1) (6,-8) rectangle ++(1,-1) (7,-8) rectangle ++(1,-1) (8,-8) rectangle ++(1,-1) (6,-9) rectangle ++(1,-1) (7,-9) rectangle ++(1,-1) (8,-9) rectangle ++(1,-1) (9,-9) rectangle ++(1,-1) (7,-10) rectangle ++(1,-1) (8,-10) rectangle ++(1,-1) (9,-10) rectangle ++(1,-1) (10,-10) rectangle ++(1,-1) (8,-11) rectangle ++(1,-1) (9,-11) rectangle ++(1,-1) (10,-11) rectangle ++(1,-1) (11,-11) rectangle ++(1,-1);
\draw[srcgray] (0,0) rectangle (12,-12);
\node[font=\scriptsize,align=center,anchor=north] at (6,-12.4) {sliding window\\($w=4$)};
\end{scope}
\begin{scope}[xshift=80mm]
\fill[fcompute!55] (0,-0) rectangle ++(1,-1) (0,-1) rectangle ++(1,-1) (1,-1) rectangle ++(1,-1) (0,-2) rectangle ++(1,-1) (1,-2) rectangle ++(1,-1) (2,-2) rectangle ++(1,-1) (0,-3) rectangle ++(1,-1) (1,-3) rectangle ++(1,-1) (2,-3) rectangle ++(1,-1) (3,-3) rectangle ++(1,-1) (0,-4) rectangle ++(1,-1) (1,-4) rectangle ++(1,-1) (2,-4) rectangle ++(1,-1) (3,-4) rectangle ++(1,-1) (4,-4) rectangle ++(1,-1) (5,-5) rectangle ++(1,-1) (5,-6) rectangle ++(1,-1) (6,-6) rectangle ++(1,-1) (5,-7) rectangle ++(1,-1) (6,-7) rectangle ++(1,-1) (7,-7) rectangle ++(1,-1) (5,-8) rectangle ++(1,-1) (6,-8) rectangle ++(1,-1) (7,-8) rectangle ++(1,-1) (8,-8) rectangle ++(1,-1) (9,-9) rectangle ++(1,-1) (9,-10) rectangle ++(1,-1) (10,-10) rectangle ++(1,-1) (9,-11) rectangle ++(1,-1) (10,-11) rectangle ++(1,-1) (11,-11) rectangle ++(1,-1);
\draw[srcgray] (0,0) rectangle (12,-12);
\node[font=\scriptsize,align=center,anchor=north] at (6,-12.4) {packed docs\\(block-diag.\ causal)};
\end{scope}
\begin{scope}[xshift=120mm]
\fill[fcompute!55] (0,-0) rectangle ++(1,-1) (1,-0) rectangle ++(1,-1) (2,-0) rectangle ++(1,-1) (3,-0) rectangle ++(1,-1) (4,-0) rectangle ++(1,-1) (0,-1) rectangle ++(1,-1) (1,-1) rectangle ++(1,-1) (2,-1) rectangle ++(1,-1) (3,-1) rectangle ++(1,-1) (4,-1) rectangle ++(1,-1) (0,-2) rectangle ++(1,-1) (1,-2) rectangle ++(1,-1) (2,-2) rectangle ++(1,-1) (3,-2) rectangle ++(1,-1) (4,-2) rectangle ++(1,-1) (0,-3) rectangle ++(1,-1) (1,-3) rectangle ++(1,-1) (2,-3) rectangle ++(1,-1) (3,-3) rectangle ++(1,-1) (4,-3) rectangle ++(1,-1) (0,-4) rectangle ++(1,-1) (1,-4) rectangle ++(1,-1) (2,-4) rectangle ++(1,-1) (3,-4) rectangle ++(1,-1) (4,-4) rectangle ++(1,-1) (0,-5) rectangle ++(1,-1) (1,-5) rectangle ++(1,-1) (2,-5) rectangle ++(1,-1) (3,-5) rectangle ++(1,-1) (4,-5) rectangle ++(1,-1) (5,-5) rectangle ++(1,-1) (0,-6) rectangle ++(1,-1) (1,-6) rectangle ++(1,-1) (2,-6) rectangle ++(1,-1) (3,-6) rectangle ++(1,-1) (4,-6) rectangle ++(1,-1) (5,-6) rectangle ++(1,-1) (6,-6) rectangle ++(1,-1) (0,-7) rectangle ++(1,-1) (1,-7) rectangle ++(1,-1) (2,-7) rectangle ++(1,-1) (3,-7) rectangle ++(1,-1) (4,-7) rectangle ++(1,-1) (5,-7) rectangle ++(1,-1) (6,-7) rectangle ++(1,-1) (7,-7) rectangle ++(1,-1) (0,-8) rectangle ++(1,-1) (1,-8) rectangle ++(1,-1) (2,-8) rectangle ++(1,-1) (3,-8) rectangle ++(1,-1) (4,-8) rectangle ++(1,-1) (5,-8) rectangle ++(1,-1) (6,-8) rectangle ++(1,-1) (7,-8) rectangle ++(1,-1) (8,-8) rectangle ++(1,-1) (0,-9) rectangle ++(1,-1) (1,-9) rectangle ++(1,-1) (2,-9) rectangle ++(1,-1) (3,-9) rectangle ++(1,-1) (4,-9) rectangle ++(1,-1) (5,-9) rectangle ++(1,-1) (6,-9) rectangle ++(1,-1) (7,-9) rectangle ++(1,-1) (8,-9) rectangle ++(1,-1) (9,-9) rectangle ++(1,-1) (0,-10) rectangle ++(1,-1) (1,-10) rectangle ++(1,-1) (2,-10) rectangle ++(1,-1) (3,-10) rectangle ++(1,-1) (4,-10) rectangle ++(1,-1) (5,-10) rectangle ++(1,-1) (6,-10) rectangle ++(1,-1) (7,-10) rectangle ++(1,-1) (8,-10) rectangle ++(1,-1) (9,-10) rectangle ++(1,-1) (10,-10) rectangle ++(1,-1) (0,-11) rectangle ++(1,-1) (1,-11) rectangle ++(1,-1) (2,-11) rectangle ++(1,-1) (3,-11) rectangle ++(1,-1) (4,-11) rectangle ++(1,-1) (5,-11) rectangle ++(1,-1) (6,-11) rectangle ++(1,-1) (7,-11) rectangle ++(1,-1) (8,-11) rectangle ++(1,-1) (9,-11) rectangle ++(1,-1) (10,-11) rectangle ++(1,-1) (11,-11) rectangle ++(1,-1);
\draw[srcgray] (0,0) rectangle (12,-12);
\node[font=\scriptsize,align=center,anchor=north] at (6,-12.4) {prefix-LM\\(prefix $=5$)};
\end{scope}
\end{tikzpicture}
\caption{Attention masks as query (row) $\times$ key (column) patterns; shaded cells are
allowed. Tiles that are entirely unshaded can be skipped by a tiled kernel.}
\label{fig:attention:patterns}
\end{figure}

\paragraph{Changing the algorithm.} Linear attention (Katharopoulos et al., 2020) replaces
$\exp(\vect{q}\cdot\vect{k})$ by a kernel feature map so that $\mat{A}\mat{V}$ can be computed as
$\phi(\mat{Q})(\phi(\mat{K})^\top\mat{V})$ in $O(Td^2)$, turning attention into a recurrence with a
fixed-size state; state-space models and their hybrids pursue the same goal. These trade exact
retrieval over the whole context for constant per-token cost.

\paragraph{Distributing the sequence.} When even $O(T)$ memory per device is too much, the sequence
itself is split across GPUs. Ring attention (Liu et al., 2023) passes key/value blocks around a ring
of devices while each device keeps its query block, using exactly the online-softmax merge of
\cref{eq:attention:online} to combine partial results. \Cref{ch:distributed} covers this context
parallelism, and \cref{ch:serving} covers PagedAttention, which changes \emph{where} keys and values
live in memory rather than which ones are attended to.

\begin{napkinmath}[Choosing a tile size]
An H100 SM has up to 228\,KB of shared memory. With $D=128$ in bf16, one $64\times128$ tile is 16\,KB.
A forward block holding $\mat{Q}_i$ ($B_r=128$ rows: 32\,KB), double-buffered $\mat{K}_j$ and
$\mat{V}_j$ ($B_c=128$: $2\times2\times 32$\,KB $=128$\,KB) uses 160\,KB, leaving headroom; the
$128\times128$ fp32 score tile lives in registers spread across the block's warps. Doubling $D$ to
256 halves the feasible tile sizes, which is one reason head dimensions above 128 are less
efficient and kernels specialise by head dimension.
\end{napkinmath}

% =====================================================================
\section{Summary}\label{sec:attention:summary}
\begin{itemize}
  \item Attention replaces the RNN's single context vector with a per-output, data-dependent
  weighted average over all inputs; self-attention applies this within one sequence and parallelises
  across positions.
  \item Scaled dot-product attention projects inputs to queries, keys and values,
  computes $\softmax(\mat{Q}\mat{K}^\top/\sqrt{d_k})\mat{V}$, and scales by $\sqrt{d_k}$ to keep the softmax
  out of saturation.
  \item Causal masking (scores set to $-\infty$ above the diagonal) prevents a position from seeing its
  future; masking after the softmax with renormalisation is equivalent. Dropout on attention weights
  is a training-only regulariser.
  \item Multi-head attention runs $H$ heads of width $d/H$; the weight-split implementation computes
  all heads with a few large batched GEMMs and an output projection.
  \item Attention costs $O(T^2d)$ FLOPs and, naively, $O(BHT^2)$ memory; at long context the memory and
  the HBM traffic of the $T\times T$ matrix, not the FLOPs, are the binding constraints.
  \item FlashAttention tiles $\mat{Q},\mat{K},\mat{V}$, uses the online softmax to normalise incrementally,
  stores one log-sum-exp per row, and recomputes probabilities in the backward pass: exact attention
  with $O(T)$ memory and far less HBM traffic.
  \item GQA/MQA, windowed and sparse masks, linear attention and ring attention modify dense attention
  to cut KV-cache size, compute, or per-device memory.
\end{itemize}

% =====================================================================
\section{Exercises}\label{sec:attention:exercises}

\begin{exercise}[hand computation]
Using the running example, compute by hand (or in five lines of PyTorch) the context vector for the
query ``chip'' with weight-free attention, and verify it against the last row of
\cref{fig:attention:heat}(a). Then compute its causal context vector.
\end{exercise}

\begin{exercise}[scaling]
Draw $\vect{q},\vect{k}\in\R^{d_k}$ with i.i.d.\ standard normal entries for $d_k\in\{16,64,256,1024\}$
and $T=512$ keys. Plot the average maximum softmax probability with and without the
$1/\sqrt{d_k}$ scale. Relate what you see to the softmax gradient
$\partial\alpha_i/\partial s_j=\alpha_i(\delta_{ij}-\alpha_j)$.
\end{exercise}

\begin{exercise}[Parameter vs.\ Linear]
Implement self-attention with raw \code{nn.Parameter} matrices used as \code{x @ W}. Show how to
copy weights between it and \cref{lst:attention:sa} so that both produce identical outputs.
\end{exercise}

\begin{exercise}[no leakage]
Prove algebraically that masking after the softmax followed by renormalisation equals masking with
$-\infty$ before the softmax. Then show numerically that, for route 1, the gradient of position 2's
output with respect to position 5's input is exactly zero.
\end{exercise}

\begin{exercise}[head bookkeeping]
GPT-2 medium has $d=1024$ and 16 heads. Give the shapes of every intermediate tensor in
\cref{lst:attention:mha} for $B=4$, $T=1024$, and the number of parameters in the attention block
(with Q/K/V biases, as in the released weights).
\end{exercise}

\begin{exercise}[FLOP crossover]
For a 32-layer model with $d=4096$, $H=32$ and SwiGLU MLP width $11008$, find the sequence length at
which attention-score FLOPs ($\mat{Q}\mat{K}^\top$ and $\mat{A}\mat{V}$, with causal masking exploited)
equal all other per-token forward FLOPs.
\end{exercise}

\begin{exercise}[online softmax algebra]
Prove by induction that \cref{eq:attention:online} yields the global $(m,\ell)$. Then show that two
partial states $(m_1,\ell_1,\vect{o}_1)$ and $(m_2,\ell_2,\vect{o}_2)$ computed over \emph{disjoint}
key sets can be merged into the state for their union, and write the merge function. Where does this
merge appear in split-KV decoding and in ring attention?
\end{exercise}

\begin{exercise}[the $\delta$ identity]
Verify numerically that $\operatorname{rowsum}(\mat{P}\odot\mat{dP})=\operatorname{rowsum}(\mat{dO}\odot\mat{O})$,
then derive the gradient of a row softmax from scratch to see why $\vect{\delta}$ appears.
\end{exercise}

\begin{exercise}[IO model]
Using \code{llmsys.flash.attention\_hbm\_bytes}, tabulate naive vs.\ tiled HBM traffic for
$T\in\{1\text{k},4\text{k},16\text{k},64\text{k}\}$ and $D\in\{64,128,256\}$ with 228\,KB of SRAM.
At what $D$ does tiling stop helping in this model? Discuss what the model leaves out.
\end{exercise}

\begin{exercise}[skipping masked tiles]
For causal attention with $T=8192$ and square tiles of 128, what fraction of tiles is fully masked,
fully visible, and on the diagonal? How much of the naive $\mat{Q}\mat{K}^\top$ work does a
tile-skipping kernel avoid?
\end{exercise}

\begin{exercise}[sliding window cache]
A model uses a sliding window of $w=4096$ in every layer. Derive the KV-cache size as a function of
generated length and compare to full attention at $T=32$k. What breaks if you want to add prefix
caching across requests?
\end{exercise}

\begin{exercise}[measurement]
Run \ident{labs/lab04_attention/bench.py} on a GPU for $T$ from 512 to 16k. Plot time and peak
memory of the explicit implementation and SDPA. Identify the length at which the explicit version
runs out of memory and compare with the napkin estimate of \cref{sec:attention:cost}.
\end{exercise}

% =====================================================================
\section{Lab 4: Attention, from dot products to IO-aware tiling}\label{sec:attention:lab}

\begin{labbox}{4}{Attention, from dot products to IO-aware tiling}
\labdir{lab04_attention}

\labsection{Goals}
Implement every attention variant of this chapter, prove to yourself (with tests) that the
fast versions compute the same function as the simple ones, and measure where the time and memory
go as the sequence grows.

\labsection{Setup}
\begin{lstlisting}[style=shell]
cd code
pytest -q labs/lab04_attention                     # grades starter.py
LAB_IMPL=solution pytest -q labs/lab04_attention   # reference
python labs/lab04_attention/bench.py [--device cuda]
\end{lstlisting}

\labsection{Tasks}
\begin{enumerate}
  \item \textbf{Weight-free attention.} \ident{simple_attention}: scores, softmax, weighted sum.
  The test checks one row against a hand-written loop.
  \item \textbf{Trainable self-attention.} \ident{SelfAttention.forward} with the $1/\sqrt{d_k}$ scale.
  \item \textbf{Causal masking.} \ident{causal_softmax} and \ident{mask_after_softmax}; they must
  agree. \ident{CausalAttention.forward} with dropout: the test perturbs future inputs and checks
  past outputs are unchanged, and that dropout is active only in \code{train()} mode.
  \item \textbf{Multi-head.} \ident{MultiHeadAttentionWrapper.forward} and
  \ident{MultiHeadAttention.forward}. Your weight-split module must produce the same output as the
  core \code{llmsys.layers.MultiHeadAttention} after \code{load\_state\_dict}, and with one head it must
  reduce to \code{out\_proj(CausalAttention(x))}.
  \item \textbf{Online softmax.} \ident{online_softmax} for arbitrary block sizes.
  \item \textbf{Tiled attention.} \ident{tiled_attention}: never allocate a $T\times T$ tensor; skip
  key tiles that are entirely in the causal future. Must match naive attention to $10^{-10}$ in
  float64 for non-divisible tile sizes.
  \item \textbf{Napkin.} \ident{score_matrix_bytes}.
\end{enumerate}

\labsection{Measurements}
\begin{itemize}
  \item Wrapper vs.\ weight-split MHA latency for GPT-2-small shapes. On a GPU expect the weight-split
  version to win clearly, by more at small batch; on a CPU the gap is smaller.
  \item Time and peak memory of explicit attention vs.\ SDPA for $T=256\ldots16$k. Expect explicit
  peak memory to grow as $T^2$ and SDPA's as $T$, and SDPA to be faster beyond a few hundred tokens.
  \item Your tiled reference is slow (Python loops); measure it only to see that its memory stays flat.
\end{itemize}

\labsection{Deliverables}
All tests passing; plots of time and memory vs.\ $T$; a paragraph comparing the measured
out-of-memory point to the napkin estimate.

\labsection{Stretch goals}
\begin{itemize}
  \item Implement the tiled \emph{backward} pass yourself and check it with
  \code{torch.autograd.gradcheck}; compare with \code{llmsys.flash.FlashAttentionFunction}.
  \item Add document masks (block-diagonal causal) to \ident{tiled_attention}, skipping tiles that
  cross document boundaries entirely.
  \item Port \ident{tiled_attention} to Triton following Appendix~\ref{app:triton}.
\end{itemize}
\end{labbox}
